% Options for packages loaded elsewhere
% Options for packages loaded elsewhere
\PassOptionsToPackage{unicode}{hyperref}
\PassOptionsToPackage{hyphens}{url}
%
\documentclass[
  ignorenonframetext,
  aspectratio=169,
]{beamer}
\newif\ifbibliography
\usepackage{pgfpages}
\setbeamertemplate{caption}[numbered]
\setbeamertemplate{caption label separator}{: }
\setbeamercolor{caption name}{fg=normal text.fg}
\beamertemplatenavigationsymbolshorizontal
% remove section numbering
\setbeamertemplate{part page}{
  \centering
  \begin{beamercolorbox}[sep=16pt,center]{part title}
    \usebeamerfont{part title}\insertpart\par
  \end{beamercolorbox}
}
\setbeamertemplate{section page}{
  \centering
  \begin{beamercolorbox}[sep=12pt,center]{section title}
    \usebeamerfont{section title}\insertsection\par
  \end{beamercolorbox}
}
\setbeamertemplate{subsection page}{
  \centering
  \begin{beamercolorbox}[sep=8pt,center]{subsection title}
    \usebeamerfont{subsection title}\insertsubsection\par
  \end{beamercolorbox}
}
% Prevent slide breaks in the middle of a paragraph
\widowpenalties 1 10000
\raggedbottom
\AtBeginPart{
  \frame{\partpage}
}
\AtBeginSection{
  \ifbibliography
  \else
    \frame{\sectionpage}
  \fi
}
\AtBeginSubsection{
  \frame{\subsectionpage}
}
\usepackage{iftex}
\ifPDFTeX
  \usepackage[T1]{fontenc}
  \usepackage[utf8]{inputenc}
  \usepackage{textcomp} % provide euro and other symbols
\else % if luatex or xetex
  \usepackage{unicode-math} % this also loads fontspec
  \defaultfontfeatures{Scale=MatchLowercase}
  \defaultfontfeatures[\rmfamily]{Ligatures=TeX,Scale=1}
\fi
\usepackage{lmodern}

\usetheme[]{metropolis}
\usecolortheme[]{dolphin}
\usefonttheme[]{professionalfonts}
\ifPDFTeX\else
  % xetex/luatex font selection
\fi
% Use upquote if available, for straight quotes in verbatim environments
\IfFileExists{upquote.sty}{\usepackage{upquote}}{}
\IfFileExists{microtype.sty}{% use microtype if available
  \usepackage[]{microtype}
  \UseMicrotypeSet[protrusion]{basicmath} % disable protrusion for tt fonts
}{}
\makeatletter
\@ifundefined{KOMAClassName}{% if non-KOMA class
  \IfFileExists{parskip.sty}{%
    \usepackage{parskip}
  }{% else
    \setlength{\parindent}{0pt}
    \setlength{\parskip}{6pt plus 2pt minus 1pt}}
}{% if KOMA class
  \KOMAoptions{parskip=half}}
\makeatother


\usepackage{longtable,booktabs,array}
\newcounter{none} % for unnumbered tables
\usepackage{calc} % for calculating minipage widths
\usepackage{caption}
% Make caption package work with longtable
\makeatletter
\def\fnum@table{\tablename~\thetable}
\makeatother
\usepackage{graphicx}
\makeatletter
\newsavebox\pandoc@box
\newcommand*\pandocbounded[1]{% scales image to fit in text height/width
  \sbox\pandoc@box{#1}%
  \Gscale@div\@tempa{\textheight}{\dimexpr\ht\pandoc@box+\dp\pandoc@box\relax}%
  \Gscale@div\@tempb{\linewidth}{\wd\pandoc@box}%
  \ifdim\@tempb\p@<\@tempa\p@\let\@tempa\@tempb\fi% select the smaller of both
  \ifdim\@tempa\p@<\p@\scalebox{\@tempa}{\usebox\pandoc@box}%
  \else\usebox{\pandoc@box}%
  \fi%
}
% Set default figure placement to htbp
\def\fps@figure{htbp}
\makeatother





\setlength{\emergencystretch}{3em} % prevent overfull lines

\providecommand{\tightlist}{%
  \setlength{\itemsep}{0pt}\setlength{\parskip}{0pt}}



 


\makeatletter
\@ifpackageloaded{tcolorbox}{}{\usepackage[skins,breakable]{tcolorbox}}
\@ifpackageloaded{fontawesome5}{}{\usepackage{fontawesome5}}
\definecolor{quarto-callout-color}{HTML}{909090}
\definecolor{quarto-callout-note-color}{HTML}{0758E5}
\definecolor{quarto-callout-important-color}{HTML}{CC1914}
\definecolor{quarto-callout-warning-color}{HTML}{EB9113}
\definecolor{quarto-callout-tip-color}{HTML}{00A047}
\definecolor{quarto-callout-caution-color}{HTML}{FC5300}
\definecolor{quarto-callout-color-frame}{HTML}{acacac}
\definecolor{quarto-callout-note-color-frame}{HTML}{4582ec}
\definecolor{quarto-callout-important-color-frame}{HTML}{d9534f}
\definecolor{quarto-callout-warning-color-frame}{HTML}{f0ad4e}
\definecolor{quarto-callout-tip-color-frame}{HTML}{02b875}
\definecolor{quarto-callout-caution-color-frame}{HTML}{fd7e14}
\makeatother
\makeatletter
\@ifpackageloaded{caption}{}{\usepackage{caption}}
\AtBeginDocument{%
\ifdefined\contentsname
  \renewcommand*\contentsname{Table of contents}
\else
  \newcommand\contentsname{Table of contents}
\fi
\ifdefined\listfigurename
  \renewcommand*\listfigurename{List of Figures}
\else
  \newcommand\listfigurename{List of Figures}
\fi
\ifdefined\listtablename
  \renewcommand*\listtablename{List of Tables}
\else
  \newcommand\listtablename{List of Tables}
\fi
\ifdefined\figurename
  \renewcommand*\figurename{Figure}
\else
  \newcommand\figurename{Figure}
\fi
\ifdefined\tablename
  \renewcommand*\tablename{Table}
\else
  \newcommand\tablename{Table}
\fi
}
\@ifpackageloaded{float}{}{\usepackage{float}}
\floatstyle{ruled}
\@ifundefined{c@chapter}{\newfloat{codelisting}{h}{lop}}{\newfloat{codelisting}{h}{lop}[chapter]}
\floatname{codelisting}{Listing}
\newcommand*\listoflistings{\listof{codelisting}{List of Listings}}
\makeatother
\makeatletter
\makeatother
\makeatletter
\@ifpackageloaded{caption}{}{\usepackage{caption}}
\@ifpackageloaded{subcaption}{}{\usepackage{subcaption}}
\makeatother

\usepackage{bookmark}
\IfFileExists{xurl.sty}{\usepackage{xurl}}{} % add URL line breaks if available
\urlstyle{same}
\hypersetup{
  pdftitle={VS Code},
  hidelinks,
  pdfcreator={LaTeX via pandoc}}


\title{\texorpdfstring{VS Code}{VS Code}}
\author{}
\date{}

\begin{document}
\frame{\titlepage}


\begin{frame}[fragile]{Install and Set Up VS Code}
\protect\phantomsection\label{sec-install-code}
You can skip this section if you already have a coding environment you
like; just set it up to work with
Julia\footnote<\value{beamerpauses}->[frame]{I assume that if this is
  the case, you know/can figure out how to configure your editor. If you
  aren't sure, post on Ed and we can find the instructions.}. Otherwise,
VS Code is as close to an officially supported editor for Julia as you
can get. We will follow
\href{https://www.julia-vscode.org/docs/dev/gettingstarted/}{this guide
for setting up VS Code with Julia}.

\begin{block}{Installing VS Code}
\protect\phantomsection\label{installing-vs-code}
You can download VS Code
\href{https://code.visualstudio.com/download}{here}; open the downloaded
file to install.
\end{block}

\begin{block}{Install the Julia Extension}
\protect\phantomsection\label{install-the-julia-extension}
\begin{enumerate}
\tightlist
\item
  Open VS Code.
\item
  Select View and click Extensions to open the Extension View. This view
  can also be found on the sidebar with the following logo:
  \pandocbounded{\includegraphics[keepaspectratio,alt={VS Code Extensions View Logo}]{../slides/images/vs-code-extensions.png}}
\item
  Search for \texttt{julia} in the search box. Click the green install
  button.
\item
  Restart VS Code once the installation is complete. It should
  automatically find your Julia installation; talk to Vivek if not.
\end{enumerate}

The Julia VS Code extension offers you some nice features. You can start
a REPL (an interactive Julia coding environment) by opening the
``Command Palette'' (View -\textgreater{} Command Palette, or
CTRL/CMD+SHIFT+P) and typing ``REPL'' to bring up ``Julia: Start REPL''.
You can also create \texttt{*.jl} files to write Julia code and execute
line by line. However, we will primarily use Jupyter notebooks in this
class, but this might be useful for testing code or for your project.
\end{block}

\begin{block}{Install the Jupyter Notebook Extension}
\protect\phantomsection\label{install-the-jupyter-notebook-extension}
The Jupyter Notebook extension allows you to export a Jupyter notebook
to PDF or to HTML and then to PDF.

\begin{tcolorbox}[enhanced jigsaw, arc=.35mm, bottomrule=.15mm, bottomtitle=1mm, breakable, colback=white, colbacktitle=quarto-callout-tip-color!10!white, colframe=quarto-callout-tip-color-frame, coltitle=black, left=2mm, leftrule=.75mm, opacityback=0, opacitybacktitle=0.6, rightrule=.15mm, title=\textcolor{quarto-callout-tip-color}{\faLightbulb}\hspace{0.5em}{PDF Export for Gradescope}, titlerule=0mm, toprule=.15mm, toptitle=1mm]

You will need to export every notebook to a PDF for submission to
Gradescope. Direct export to PDF requires a LaTeX installation. If you
would like to go this route, please look at
\href{https://www.latex-project.org/get/}{the LaTeX installation
instructions for your operating system}.

Otherwise, \href{homework.qmd}{exporting to HTML and then using your
browser to save the resulting page to a PDF} is a perfect solution.

\end{tcolorbox}

Follow the same instructions as above, but search for \texttt{jupyter}
and install the Jupyter extension. Restart VS Code.
\end{block}
\end{frame}

\begin{frame}{Using GitHub with VS Code}
\protect\phantomsection\label{using-github-with-vs-code}
\href{https://youtu.be/z5jZ9lrSpqk?t=115&si=PcpNxjmF1ByTUdJq}{This
Youtube video} provides an overview of how to use GitHub within VS Code
(with chapters for specific steps) and may be worth watching.
\end{frame}




\end{document}
